% !PRINTED 651
\weberpage{651}

If, then, $\xi$ is a number of the field $\Omega$ that belongs to the
group $\mathfrak{A}$, then $\xi$ is a primitive number of the field
$R(\xi)$ of degree $b$, and this field also contains every number $\xi'$
that admits the substitutions of the group $\mathfrak{A}'$. If now
$\mathfrak{A}, \mathfrak{B}$ and $\mathfrak{A}', \mathfrak{B}'$ are proper divisors of $\mathfrak{G}$, then $R(\xi)$ is different from
$\Omega$, and $R(\xi)$ is not extended by the adjunction of $\xi'$.
Hence $\Omega$ also cannot be composed from $R(\xi)$ and $R(\xi')$,
by which 3. is proved.

\begin{center}
\S.\,180.

\gesperrt{The Resolvents.}
\end{center}

After what has now been proved, we may restrict ourselves, in
the further considerations that have the proof of Theorem I. as
their goal, to those \gesperrt{Abelian} fields whose degree is a prime power
and whose group is cyclic. For this reason it also sufficed for
the intended application that in the preceding chapter we confined
ourselves to the consideration of the cyclotomic fields $\Omega_m$,
in which $m$ was a prime power. We retain, as far as possible,
the notation used there, and set
\weq{1}{m = q^{\varkappa},}
where $q$ is a prime and $\varkappa$ a positive exponent, and in the
case that $q = 2$ we assume $\varkappa > 1$. By $n$ we denote every
number not divisible by $q$. Here $r$ is a primitive $m$th root of
unity and $\Omega_m$ is the field of the rational functions of $r$.

Now let $\Omega = R(x)$ be a simple \gesperrt{Abelian} field of degree $m$,
so that $x$ is a root of a rational irreducible cyclic equation of
degree $m$. If we denote the roots of this equation, in a suitable
order, by
\weq{2}{x_0, x_1, x_2, \ldots, x_{m-1},}
and take $x_m = x_0$, $x_{m+1} = x_1, \ldots$, then
\[
x_1 = \Phi(x_0), \quad x_2 = \Phi(x_1), \quad \ldots, \quad x_0 = \Phi(x_{m-1}),
\]
where $\Phi$ denotes an integral function with rational coefficients,
and the group of the field $\Omega$ consists of the powers of the
substitution $(x_0, x_1)$ or of the cyclic permutations
%CONT>>
